% ============================================================
%  capitoli/frazioni.tex — Le frazioni
%  Matemagica — Matematica per le scuole medie
%  Autori: Francesco & Emma Piraneo Giuliano
%  Licenza: CC BY-NC-SA 4.0
% ============================================================

\chapter{Le frazioni}

% ------------------------------------------------------------
\section*{Obiettivi di apprendimento}
% ------------------------------------------------------------

Al termine di questo capitolo sarai in grado di:

\begin{itemize}[leftmargin=1.5em]
    \item riconoscere e leggere una frazione;
    \item identificare il \textbf{numeratore} e il \textbf{denominatore};
    \item rappresentare una frazione come parte di un intero;
    \item utilizzare la frazione come \textbf{operatore}, cioè per calcolare una parte di una quantità.
\end{itemize}

% ------------------------------------------------------------
\section{Frazioni nella vita di tutti i giorni}
% ------------------------------------------------------------

Osserva queste frasi, tratte dalla vita quotidiana:

\begin{itemize}[leftmargin=1.5em]
    \item «Mi dia mezzo litro di latte e un chilo e mezzo di pane, per favore.»
    \item «Mancano ancora tre quarti d'ora alla fine della lezione.»
    \item «La mia squadra ha vinto l'incontro dopo aver dominato per i due terzi della partita.»
    \item «Nel negozio all'angolo vengono praticati sconti del trenta percento.»
\end{itemize}

In tutti questi casi si parla di \textbf{frazioni}: simboli che indicano una parte di qualcosa.
Anche le percentuali sono frazioni: il $30\%$ significa, letteralmente, \emph{trenta centesimi},
cioè $\dfrac{30}{100}$.

% ------------------------------------------------------------
\section{La frazione come parte di un intero}
% ------------------------------------------------------------

\begin{definizione}
    Si chiama \textbf{frazione} il simbolo $\dfrac{a}{b}$, con $b \neq 0$,
    dove $a$ e $b$ sono numeri interi.
    Una frazione rappresenta una o più parti uguali in cui è stato diviso un intero.
\end{definizione}

Il modo più immediato per capire le frazioni è immaginare una torta.

\begin{figure}[htb]
    \centering
    \begin{tikzpicture}
        % --- torta intera suddivisa in 8 spicchi, 3 colorati ---
        \def\R{2.2}
        % Sfondo grigio chiaro = spicchi non colorati
        \foreach \i in {0,1,2,3,4,5,6,7} {
                \fill[gray!20] (0,0) -- ({(\i)*45}:\R) arc[start angle={(\i)*45},
                        end angle={(\i+1)*45}, radius=\R] -- cycle;
            }
        % Spicchi colorati (blu chiaro) = i primi 3
        \foreach \i in {0,1,2} {
                \fill[blu!40] (0,0) -- ({(\i)*45}:\R) arc[start angle={(\i)*45},
                        end angle={(\i+1)*45}, radius=\R] -- cycle;
            }
        % Bordi
        \draw[thick, blu] (0,0) circle (\R);
        \foreach \i in {0,1,2,3,4,5,6,7} {
                \draw[thick, blu] (0,0) -- ({(\i)*45}:\R);
            }
        % Etichetta
        \node[below] at (0,-\R-0.4) {$\dfrac{3}{8}$ della torta};
    \end{tikzpicture}
    \caption{Una torta divisa in 8 parti uguali: 3 parti su 8 sono colorate.}
\end{figure}

La torta è stata divisa in \textbf{8 parti uguali}. Ne abbiamo prese \textbf{3}.
Scriviamo quindi $\dfrac{3}{8}$ e leggiamo: \emph{«tre ottavi»}.

\begin{definizione}
    Nella frazione $\dfrac{a}{b}$:
    \begin{itemize}[leftmargin=1.5em]
        \item $a$ si chiama \textbf{numeratore} e indica \emph{quante parti} consideriamo;
        \item $b$ si chiama \textbf{denominatore} e indica \emph{in quante parti uguali}
              è stato diviso l'intero;
        \item la linea orizzontale tra i due numeri si chiama \textbf{linea di frazione}.
    \end{itemize}
\end{definizione}

\begin{figure}[htb]
    \centering
    \begin{tikzpicture}[>=stealth, thick]
        % Frazione grande al centro
        \node[font=\Large] (num) at (0, 0.7)  {$a$};
        \node[font=\Large] (bar) at (0, 0)    {};
        \draw[blu, very thick] (-0.4, 0) -- (0.4, 0);
        \node[font=\Large] (den) at (0,-0.7)  {$b$};
        % Freccia numeratore
        \draw[blu, ->] (1.8, 0.7) node[right, font=\small\bfseries] {numeratore} -- (0.25, 0.7);
        % Freccia linea di frazione
        \draw[blu, ->] (-2.4, 0) node[left, font=\small\bfseries] {linea di frazione} -- (-0.45, 0);
        % Freccia denominatore
        \draw[blu, ->] (1.8,-0.7) node[right, font=\small\bfseries] {denominatore} -- (0.25,-0.7);
    \end{tikzpicture}
    \caption{Le parti di una frazione.}
\end{figure}

\begin{regola}
    Per \textbf{leggere} una frazione si pronuncia prima il numeratore
    (come numero cardinale) e poi il denominatore (come numero ordinale):
    \begin{itemize}[leftmargin=1.5em, itemsep=2pt]
        \item $\dfrac{1}{2}$ si legge \emph{un mezzo}
        \item $\dfrac{3}{4}$ si legge \emph{tre quarti}
        \item $\dfrac{5}{8}$ si legge \emph{cinque ottavi}
        \item $\dfrac{7}{10}$ si legge \emph{sette decimi}
    \end{itemize}
    Per denominatori maggiori di 10 si aggiunge il suffisso \emph{-esimi}:
    $\dfrac{3}{11}$ si legge \emph{tre undicesimi}.
\end{regola}

% --- Esempio visivo: torta + rettangolo ----------------------

\begin{esempio}
    \textbf{Rappresentazioni della stessa frazione in forme diverse.}

    \medskip
    La frazione $\dfrac{2}{5}$ può essere rappresentata sia con una figura circolare
    sia con un rettangolo: la forma non cambia il significato.

    \medskip
    \begin{center}
        \begin{tikzpicture}
            % --- Cerchio: 2/5 colorati ---
            \def\R{1.6}
            \foreach \i in {0,1,2,3,4} {
                    \fill[gray!20] (0,0) -- ({90 + (\i)*72}:\R)
                    arc[start angle={90+(\i)*72}, end angle={90+(\i+1)*72}, radius=\R] -- cycle;
                }
            \foreach \i in {0,1} {
                    \fill[blu!40] (0,0) -- ({90 + (\i)*72}:\R)
                    arc[start angle={90+(\i)*72}, end angle={90+(\i+1)*72}, radius=\R] -- cycle;
                }
            \draw[thick, blu] (0,0) circle (\R);
            \foreach \i in {0,1,2,3,4} {
                    \draw[thick, blu] (0,0) -- ({90+(\i)*72}:\R);
                }
            \node[below] at (0,-\R-0.35) {$\dfrac{2}{5}$};
        \end{tikzpicture}
        \hspace{2cm}
        \begin{tikzpicture}
            % --- Rettangolo: 5 colonne, 2 colorate ---
            \def\W{0.9}  % larghezza di ogni cella
            \def\H{1.6}  % altezza
            \foreach \i in {0,1,2,3,4} {
                    \fill[gray!20] ({\i*\W}, 0) rectangle ({(\i+1)*\W}, \H);
                }
            \foreach \i in {0,1} {
                    \fill[blu!40] ({\i*\W}, 0) rectangle ({(\i+1)*\W}, \H);
                }
            \draw[thick, blu] (0,0) rectangle ({5*\W}, \H);
            \foreach \i in {1,2,3,4} {
                    \draw[thick, blu] ({\i*\W},0) -- ({\i*\W},\H);
                }
            \node[below] at ({5*\W/2}, -0.35) {$\dfrac{2}{5}$};
        \end{tikzpicture}
    \end{center}

    \medskip
    In entrambe le figure l'intero è diviso in \textbf{5 parti uguali}
    e \textbf{2} di esse sono colorate: la frazione è $\dfrac{2}{5}$ in ogni caso.
\end{esempio}

\begin{attenzione}
    \textbf{Il denominatore non può essere zero.}

    Scrivere $\dfrac{3}{0}$ non ha senso matematico: non è possibile dividere
    qualcosa in \emph{zero} parti uguali.
    Per questo motivo la condizione $b \neq 0$ è sempre obbligatoria.
\end{attenzione}

% ------------------------------------------------------------
\section{La frazione come operatore}
% ------------------------------------------------------------

La frazione non serve solo a descrivere una parte di una figura:
può anche essere applicata a una \textbf{quantità concreta} per calcolarne una parte.
In questo caso si dice che la frazione agisce da \textbf{operatore}.

\begin{regola}
    Per calcolare $\dfrac{a}{b}$ di una quantità $Q$:
    \begin{enumerate}[leftmargin=1.5em, itemsep=2pt]
        \item si divide $Q$ per il denominatore $b$ (si ottiene il valore di \emph{una} parte);
        \item si moltiplica il risultato per il numeratore $a$ (si ottengono le $a$ parti cercate).
    \end{enumerate}
    \[
        \frac{a}{b} \text{ di } Q \;=\; (Q \div b) \times a
    \]
\end{regola}

\begin{esempio}
    \textbf{Luigi ha percorso $\dfrac{3}{5}$ di una pista lunga 350\,m.
        Quanti metri ha percorso?}

    \medskip
    Passo 1 — divido la pista in 5 parti uguali:
    \[
        350 \div 5 = 70 \text{ m per ogni parte}
    \]

    Passo 2 — prendo 3 di quelle parti:
    \[
        70 \times 3 = 210 \text{ m}
    \]

    \medskip
    \textbf{Rappresentazione grafica (barra frazionaria):}

    \medskip
    \begin{center}
        \begin{tikzpicture}
            \def\TotW{7.5}   % larghezza totale della barra
            \def\H{0.7}      % altezza
            \def\W{\TotW/5}  % larghezza di ogni parte
            % 5 celle
            \foreach \i in {0,1,2,3,4} {
                    \fill[gray!20] ({\i*\W},0) rectangle ({(\i+1)*\W},\H);
                }
            % 3 celle colorate
            \foreach \i in {0,1,2} {
                    \fill[blu!40] ({\i*\W},0) rectangle ({(\i+1)*\W},\H);
                }
            % bordi
            \draw[thick,blu] (0,0) rectangle (\TotW,\H);
            \foreach \i in {1,2,3,4} {
                    \draw[thick,blu] ({\i*\W},0) -- ({\i*\W},\H);
                }
            % etichette sotto ogni cella
            \foreach \i in {0,1,2,3,4} {
                    \node[below, font=\small] at ({(\i+0.5)*\W}, -0.05) {70\,m};
                }
            % parentesi graffa per le 3 parti colorate
            \draw[blu, thick, decorate,
                decoration={brace, amplitude=6pt, raise=4pt}]
            (0, \H) -- ({3*\W}, \H)
            node[midway, above=10pt, font=\small\bfseries] {$\frac{3}{5}$ di 350\,m $= 210$\,m};
            % etichetta totale sotto
            \node[below=1.1cm, font=\small] at ({\TotW/2},0) {350\,m totali};
        \end{tikzpicture}
    \end{center}

    \medskip
    Luigi ha percorso \textbf{210 m}.
\end{esempio}

\begin{attenzione}
    \textbf{Errore comune: confondere numeratore e denominatore nell'operatore.}

    Alcuni allievi, per calcolare $\dfrac{3}{5}$ di 350, dividono per 3 e moltiplicano per 5.

    \medskip
    \begin{tabular}{ll}
        Procedimento errato:   & $350 \div 3 \times 5 \approx 583$ \\[4pt]
        Procedimento corretto: & $350 \div 5 \times 3 = 210$
        \quad\textit{(corretto)}
    \end{tabular}

    \medskip
    Ricorda: si divide \textbf{sempre} per il denominatore (il numero \emph{in basso}).
\end{attenzione}

% ------------------------------------------------------------
\section*{Lo sapevi?}
% ------------------------------------------------------------

\begin{sapevi}
    \textbf{Le frazioni nell'antico Egitto e il Papiro di Rhind}

    Le frazioni sono uno degli strumenti matematici più antichi che conosciamo.
    Gli antichi Egizi le usavano già oltre \textbf{3\,500 anni fa}, come testimoniano
    diversi papiri conservati fino ai giorni nostri.

    Il più importante è il \textbf{Papiro di Rhind} (o Papiro di Ahmes),
    scritto attorno al 1650 a.C. dal matematico e scriba Ahmes.
    Si tratta di una raccolta di 87 problemi matematici, molti dei quali riguardano
    proprio le frazioni.

    Gli Egizi avevano però un'abitudine diversa dalla nostra: usavano quasi
    esclusivamente le cosiddette \textbf{frazioni unitarie}, cioè frazioni con
    numeratore uguale a 1, come $\dfrac{1}{2}$, $\dfrac{1}{3}$, $\dfrac{1}{7}$.
    Per esprimere frazioni più grandi, le \emph{sommavano} tra loro.
    Ad esempio, quella che noi scriviamo $\dfrac{2}{5}$ gli Egizi la scrivevano come
    $\dfrac{1}{3} + \dfrac{1}{15}$.

    Un risultato notevole, vero? Dimostra che già migliaia di anni fa gli esseri umani
    avevano sviluppato strumenti precisi per ragionare sulle parti di un intero.
\end{sapevi}

\begin{sapevi}
    \textbf{Chiavi inglesi: millimetri contro frazioni di pollice}

    Se hai mai messo mano a una cassetta degli attrezzi, potresti aver notato
    che le chiavi meccaniche non sono tutte uguali: alcune riportano misure
    in \textbf{millimetri} (sistema metrico), altre in \textbf{frazioni di pollice}
    (sistema imperiale, usato negli Stati Uniti).

    Nel sistema metrico confrontare le chiavi è immediato: una chiave da 13\,mm
    è più grande di una da 11\,mm — basta leggere i numeri interi.

    Nel sistema imperiale, invece, le stesse misure vengono espresse come frazioni
    di pollice: $\dfrac{7}{16}$\,in, $\dfrac{1}{2}$\,in, $\dfrac{9}{16}$\,in,
    $\dfrac{5}{8}$\,in, \ldots\ e capire subito quale chiave sia più grande
    richiede di \textbf{confrontare frazioni con denominatori diversi}:
    non è affatto immediato!

    Ad esempio: $\dfrac{9}{16}$ o $\dfrac{5}{8}$? Qual è la chiave più grande?

    Per rispondere bisogna trovare un denominatore comune: $\dfrac{5}{8} = \dfrac{10}{16}$,
    quindi $\dfrac{10}{16} > \dfrac{9}{16}$: la chiave da $\dfrac{5}{8}$\,in è più grande.

    Ecco un caso concreto in cui il sistema metrico — basato su numeri interi — rende
    la vita molto più semplice delle frazioni. Ma capire come funziona il confronto
    tra frazioni è fondamentale: lo affronteremo in un capitolo successivo.
\end{sapevi}

% ------------------------------------------------------------
\section*{Riepilogo del capitolo}
% ------------------------------------------------------------

\begin{tcolorbox}[
        enhanced,
        breakable,
        colback  = blu!5,
        colframe = blu,
        title    = {Riepilogo — Le frazioni},
        fonttitle = \bfseries,
        arc      = 4pt,
    ]
    \begin{itemize}[leftmargin=1.5em, itemsep=4pt]
        \item Una \textbf{frazione} $\dfrac{a}{b}$ (con $b \neq 0$) rappresenta
              $a$ parti di un intero diviso in $b$ parti uguali.
        \item $a$ si chiama \textbf{numeratore}, $b$ si chiama \textbf{denominatore}.
        \item Il denominatore non può mai essere \textbf{zero}.
        \item Una frazione può agire da \textbf{operatore}: per calcolare $\dfrac{a}{b}$
              di una quantità $Q$, si calcola $(Q \div b) \times a$.
        \item Si divide sempre per il \textbf{denominatore} (numero in basso),
              poi si moltiplica per il \textbf{numeratore} (numero in alto).
    \end{itemize}
\end{tcolorbox}